% ======================================================================
% APPENDED SUMMARIES (English + Dutch)
% Place this file at the end of your thesis (after the bibliography)
% ======================================================================

\cleardoublepage
\chapter*{English Summary}
\addcontentsline{toc}{chapter}{English Summary}

This thesis explored the complex relationship between margin adjustment,
template design, and the cognitive endurance of novice researchers.
Across multiple chapters, it examined how formatting constraints,
iterative LaTeX compilation errors, and aesthetic perfectionism interact
to shape the emotional trajectory of early-career academics.

\vspace{1em}
The main findings indicate that:
\begin{itemize}
  \item The perceived control over layout strongly predicts motivation during the writing process.
  \item Periodic frustration spikes (\emph{e.g.}, missing \texttt{\textbackslash end\{document\}}) correlate with higher eventual mastery.
  \item Collaboration through shared Overleaf projects increases both productivity and collective despair tolerance.
  \item Visual alignment consistency contributes disproportionately to subjective satisfaction with the final product.
\end{itemize}

In conclusion, the study underscores that template development is not
merely a technical activity but a deeply emotional, iterative process.
Future work should investigate margin optimization strategies and the
psychological recovery period following prolonged exposure to LaTeX warnings.

% ======================================================================
% SAMENVATTING IN HET NEDERLANDS
% ======================================================================

\cleardoublepage
\chapter*{Nederlandse Samenvatting}
\addcontentsline{toc}{chapter}{Nederlandse Samenvatting}

Dit proefschrift onderzoekt de complexe relatie tussen marge-instellingen,
sjabloonontwerp en de cognitieve belasting van beginnende onderzoekers.
In verschillende hoofdstukken is gekeken hoe opmaakbeperkingen,
terugkerende LaTeX-foutmeldingen en esthetisch perfectionisme samen
de emotionele loopbaan van promovendi beïnvloeden.

\vspace{1em}
De belangrijkste bevindingen zijn:
\begin{itemize}
  \item De ervaren controle over de opmaak voorspelt de motivatie tijdens het schrijven.
  \item Fases van frustratie (zoals een ontbrekende \texttt{\textbackslash end\{document\}}) hangen samen met latere beheersing van LaTeX.
  \item Samenwerking via gedeelde Overleaf-projecten verhoogt zowel de productiviteit als de gezamenlijke frustratietolerantie.
  \item Consistente uitlijning van tekst en tabellen draagt sterk bij aan tevredenheid over het eindresultaat.
\end{itemize}

Samenvattend laat dit onderzoek zien dat sjabloonontwikkeling niet enkel
een technische, maar ook een emotionele en iteratieve activiteit is.
Toekomstig onderzoek kan zich richten op efficiëntere margestrategieën
en het herstelproces na langdurige blootstelling aan LaTeX-waarschuwingen.